\documentclass[10pt,a4paper]{amsart}
\usepackage{amsmath,amssymb,amsthm,graphicx}
\RequirePackage{amssymb,amsmath,amsthm,bm}
\RequirePackage{thmtools}
\RequirePackage[capitalize]{cleveref}
\providecommand{\tightlist}{%
  \setlength{\itemsep}{0pt}\setlength{\parskip}{0pt}}

%%%%% varphi instead of phi
\renewcommand{\phi}{\varphi}

\mathchardef\mhyph="2D 
\providecommand{\A}{\mathcal{A}}
\providecommand{\C}{\mathcal{C}}
\providecommand{\L}{\mathcal{L}}
\providecommand{\B}{\mathcal{B}}
\providecommand{\M}{\mathcal{M}}
\renewcommand{\S}{\mathcal{S}}
\providecommand{\ol}[1]{\overline{#1}}
\providecommand{\mbf}[1]{\mathbf{#1}}
\providecommand{\mf}[1]{\mathfrak{#1}}
\providecommand{\mrm}[1]{\mathrm{#1}}
\providecommand{\mc}[1]{\mathcal{#1}}
\providecommand{\ra}{\rightarrow}
\providecommand{\Ra}{\Rightarrow}
\providecommand{\La}{\Leftarrow}
\providecommand{\lr}{\leftrightarrow}
\providecommand{\hf}{\mathbb{HF}}
\providecommand{\spec}{DgSp}
\providecommand{\bSigma}{\bm\Sigma}
\providecommand{\bDelta}{\bm\Delta}
\providecommand{\bPi}{\bm\Pi}
\providecommand{\Pinf}[1]{\Pi^{\mathrm{in}}_{#1}}
\providecommand{\Sinf}[1]{\Sigma^{\mathrm{in}}_{#1}}
\providecommand{\dSinf}[1]{d\mhyph\Sigma^{\mathrm{in}}_{#1}}
\providecommand{\Dinf}[1]{\Delta^{\mathrm{in}}_{#1}}
\providecommand{\bp}{\bar p}
\providecommand{\bq}{\bar q}
\providecommand{\ba}{\bar a}
\providecommand{\bb}{\bar b}
\providecommand{\bc}{\bar c}
\providecommand{\bd}{\bar d}
\providecommand{\bx}{\bar x}
\providecommand{\by}{\bar y}
\providecommand{\bz}{\bar z}
\providecommand{\LO}{\mathrm{LO}}
\declaretheorem[name={Theorem},numberwithin=section]{theorem}
\declaretheorem[name={Proposition}, sibling=theorem]{proposition}
\declaretheorem[name={Lemma}, sibling=theorem]{lemma}
\declaretheorem[name={Fact}, sibling=theorem]{fact}
\declaretheorem[name={Claim}, numberwithin=theorem]{claim}
\declaretheorem[name={Definition},numberwithin=section,sibling=theorem, style=definition]{definition}
\declaretheorem[name={Remark}, numbered=no,style=definition]{remark}
\declaretheorem[name={Corollary}, sibling=theorem]{corollary}
\declaretheorem[name={Question},sibling=theorem]{question}
\newcommand{\T}{{\mathcal{T}}}
\newcommand{\SR}{SR}
\newcommand{\pSR}{pSR}
\newcommand{\Q}{\mathbb{Q}}
\newcommand{\Z}{\mathbb{Z}}
\newcommand{\K}{\mathcal{K}}
\newcommand{\N}{\mathcal{N}}
\renewcommand{\L}{\mathcal{L}}
\providecommand{\Ginf}[1]{\Gamma^{\mathrm{in}}_{#1}}
\title{Linear orders of exact Scott sentence complexity Sigma-in\_(lambda+3)}
\author{David Gonzalez, Dino Rossegger}
\date{arxiv:2305.07126v1; CC BY 4.0}
\begin{document}
\maketitle
\noindent\textit{Editorial note.} All ordinals in this Scott-analysis problem are countable, and limit ordinal means nonzero limit. L times Q replaces each rational point with a copy of L, as defined in1017-1020.\par\medskip
\noindent\textit{Editorial note.} The author proof at1556 calls z the second-copy initial element, while the displayed sec(x,z) formula and its subsequent use1557 specify x as that element and z as the initial element of the first copy. The formula is preserved; this variable-description slip is explicitly recorded for independent review.\par\medskip
\noindent\textit{Editorial note.} The upper-bound passage uses the standard named-parameter Scott-rank criterion stated among the inputs. The lower-bound comparison is with A+B; it is not isomorphic to A+L2+B because the latter has a least doubled-ordinal block and the former does not.\par\medskip
\noindent\textit{Editorial note.} Native mathematical macros and theorem declarations are selected from the author dino.sty. Its unrelated draft, bibliography-backend and page-layout options are not loaded. Bibliographic metadata from the supplied flattened.bbl is formatted as a standard bibliography; author mathematical spans are unchanged.\par\medskip
\noindent\textit{Editorial scope.} One exact Scott-sentence complexity construction. The complete bespoke shuffle-sum comparison and author main proof are solution obligations. Classical well-order Scott analysis and the general orbit-to-Scott-sentence criterion are explicit inputs. The previously published Z-power lemma is not re-counted or required.\par\medskip
\noindent\textit{Allowed standard prerequisites (editorial).}\begin{itemize}
\item Infinitary Scott analysis and back-and-forth preservation. A Sigma-in\_alpha formula is an infinitary countable disjunction of existential blocks over lower-level Pi formulas, with the dual definition for Pi. The back-and-forth preorder matches those formula classes. If all tuple orbits are Sigma-in\_alpha definable over a fixed finite tuple, a Sigma-in\_(alpha+2) Scott sentence exists. A nonisomorphic K with L<=\_beta K rules out a Pi-in\_beta Scott sentence of L.. Source45-86 and the stated standard background173-184.. Translate the actual orbit upper bounds and constructed nonisomorphic back-and-forth witness into exact sentence complexity.
\item Classical Scott analysis of well orders and finite condensation. For a nonzero countable limit ordinal lambda, omega\textasciicircum{}lambda <=\_(lambda+1) omega\textasciicircum{}lambda*2. For each gamma<omega\textasciicircum{}lambda its interval isomorphism type has infinitary complexity below lambda. The finite-condensation relations start with equality, take finite distance in the quotient at successors, and take unions at limits; relation sim\_delta is Sigma-in\_(2*delta) definable.. Source1314-1333,1381-1391 and1530-1544; Ash1986, Ash-Knight2000 and Alvir-Rossegger2020.. Only the well-order building blocks and definability of block relations are allowed. The shuffle-sum construction, orbit definitions and two-round back-and-forth argument are not assumed.
\end{itemize}\medskip
In this article we give a comprehensive Scott analysis of the class of linear
orderings using three definitions of rank. Given a structure $\A$ the
\emph{(parameterless) Scott rank} $\SR(\A)$, is the least $\alpha$, such that
the automorphism orbits of all tuples in $\A$ are $\Sinf{\alpha}$ definable; the
\emph{parameterized Scott rank} $\pSR(\A)$, is the least ordinal $\alpha$ such
that there is a parameter $\bp\in A^{<\omega}$ such that $SR((\A,\bp))=\alpha$;
and the \emph{Scott sentence complexity} of $\A$, $SSC(\A)$ is the least of the
complexities $\Pinf{\alpha}$, $\Sinf{\alpha}$, $\dSinf{\alpha}$ such that $\A$
has a Scott sentence of one of these complexities. While the latter is not an

\subsection{Preliminaries}
We assume that the reader is familiar with the basic techniques of Scott
analysis and with linear orderings in general. For background we refer to
Montalb\'an's upcoming book~\cite{montalban2021} and the standard reference by
Ash and Knight~\cite{ash2000}. Let us highlight the following lemma
which we will use as basic tools in our proofs without reference.
\begin{lemma}[{\cite[Lemma 15.8]{ash2000}}]\label{lem:bfandpartitions}  Suppose $\A,\B\in\LO$. Then $\A\leq_1 \B$ if and only if $\A$ is infinite or at least as large as $\B$. For $\alpha>1$, $\A\leq_\alpha \B$ if and only if for every $1\leq\beta<\alpha$ and every partition of $\B$ into intervals $\B_0,\dots,\B_n$, with endpoints in $\B$, there is a partition of $\A$ into intervals $\A_1,\dots,\A_n$ with endpoints in $\A$, such that $\B_i\leq_\beta \A_i$.
\end{lemma}
Notice that this lemma particularly implies that $(\A,\ba)\leq_\alpha (\B,\bb)$
if and only if $(-\infty, a_0)\leq_\alpha (-\infty, b_0)$,
$(a_{|\ba-1|},\infty)\leq_\alpha (b_{|\bb-1|},\infty)$ and for all $0\leq
i<|\ba-1|$ $(a_i,a_{i+1})\leq_\alpha (b_i,b_{i+1})$~\cite[Lemma 15.7]{ash2000}.

 \section{Possible Scott sentence complexities of linear orderings}\label{sec:ssc}
We now exhibit possible Scott sentence complexities of linear orderings. Many of
our examples are constructed by producing shuffle sums of linear orderings.
Recall that given a set $S$ of countable linear orderings, the \emph{shuffle
sum} $Sh(S)$ of $S$ is obtained by partitioning $\Q$ into $|S|$ mutually dense
sets $K_i$ and replacing each point in $K_i$ with a copy of $S_i$.
Note that a linear order of the form $L\cdot\Q$ is a special case of a shuffle
sum. We begin by proving a couple of lemmas about the back-and-forth relations
of shuffle sums.

\begin{theorem}\label{thm:Slim+3}
For any limit ordinal $\lambda$, there is a linear order of Scott sentence
complexity $\Sinf{\lambda+3}.$
\end{theorem}

\begin{lemma}\label{lem:shufflesum}
Let $\A$ be a shuffle sum and $\B = \L\cdot\Q$ for some ordering $\L$. If
$\A\leq_\alpha \B$ then $\A+\L+\B\leq_{\alpha+2}\A+\B$.
\end{lemma}
\begin{proof}
To show this, we describe a winning strategy for the $\exists$-player in the
$\alpha+2$-back-and-forth game. 
On the first turn, the $\exists$-player copies the selected tuple of the
$\forall$-player using an isomorphism of the initial $\A$s and the final $\B$s. 
The only possible non-isomorphic interval is the one between the smallest
element $b$ in $\B$ and the largest element $a$ in $\A$. 
In other words, we need to show that 
\[\A_{>a}+\L+\B_{<b}\geq_{\alpha+1} \A_{>a}+\B_{<b}.\]
Because $\A$ and $\B$ are shuffle sums, $\A_{>a}\cong N+\A$ and $\B_{<b}\cong
\B+K$ for some linear orderings $N$ and $K$. 
This means that we only need to show that
$\A+\L+\B\geq_{\alpha+1}\A+\B$
if we play isomorphically on $N$ and $K$.

On the next turn, the $\forall$-player plays on $\A+\L+\B$, and the
$\exists$-player responds on $\A+\B$ using an isomorphism between the initial
$\A$s and an isomorphism of $\L+\B$ to a final segment of $\B$.
Assume without loss of generality, that $c$ is the smallest element in $\L$
played by the $\forall$-player in the distinguished copy of $\L$ within $\A+\L+\B$.
The only possible non-isomorphic interval is between $c$ and the
largest element $a$ selected within $\A$.
In other words, we need to show that 
\[\A_{>a}+\L_{<c}\leq_{\alpha} \A_{>a}+\B+\L_{<c}.\]
As $\A_{>a}\cong M+\A$ and the $\exists$-player can play the isomorphism between
the copies of $\L_{<c}$ and $M$, it is sufficient to show that $\A\leq_{\alpha}\A+\B$.
As $\A$ is a shuffle sum $\A\cong \A+\A$.
Therefore, it is enough to show that $\A\leq_\alpha \B$, as desired.
\end{proof}

\begin{corollary}\label{lower}
Let $\A$ be a shuffle sum and $\B = \L\cdot\Q$ for some ordering $\L$. If $\A+\L+\B\not\cong \A+\B$ and
$\A\leq_\alpha \B$, then $\A+\L+\B$ has no $\Pinf{\alpha+2}$ Scott sentence.
\end{corollary}

\begin{proof}
We let
$\L_1:=1+\Z+\omega^\lambda$ and $\L_2:=1+\Z+\omega^\lambda\cdot2$.
One important initial observation is that $\L_1\leq_{\lambda+1}\L_2$ as $\omega^\lambda\leq_{\lambda+1}\omega^\lambda\cdot2$.
As alluded to above, we take $\A=\L_1\cdot \Q$ and $\B=\L_2\cdot\Q$.
Finally we let $\L=\A+\L_2+\B$ and claim that $\L$ has the desired complexity
Let us begin with showing that $\SR(\A)\leq\lambda$. We define several helping predicates from which the claim will easily follow.
\begin{enumerate}
    \item Let $pt(x)$ if $x$ is in a block of size $1$. 
    \item Let $Z(x)$ say $\exists z<x ~pt(z)\land \forall y~z<y<x \to
        \bigvee_{n} S^n(y)=x.$ Note that $Z(x)$ holds if and only if $x$ is inside one of the copies of $\Z$.
    \item Let $init(x)$ if $\exists z<x ~ pt(z) \land (z,x)\cong\Z.$ Note that
        $init(x)$ if and only if $x$ is the first element in a copy of $\omega^\lambda$.
    \item Let $\gamma(x)$ if $\exists z<x ~init(z)\land [z,x)\cong\gamma$ for
        any $\gamma<\omega^\lambda$. Note that this definition is of complexity
        lower than $\Sinf\lambda$ as $\lambda$ is a limit ordinal. It describes
        the points that are of rank $\gamma$ within a copy of $\omega^{\lambda}$.
\end{enumerate}
Using these predicates one can define the automorphism orbits of any tuples in
$\A$ using formulas of complexity less than $\Sinf{\lambda+1}$.

All of these predicates have the same meaning within $\B$, however they leave
points in the second copy of $\omega^\lambda$ undefined. But it is still the case that $\SR(B)\leq\lambda+1$. To see this we need to define the additional relations
\begin{enumerate}
    \item Let 
        \[sec(x,z)\iff init(z) \land z<x \land \lnot\exists y~ \left(z<y<x\land
            pt(y)\right) \land \bigwedge_{\delta<\lambda}\forall y<x~ y\not\sim_\delta
        x.\]
        This $\Pinf\lambda$ predicate holds only if $z$ is the initial element in the second copy of $\omega^{\lambda}$ above the copy that $z$ is initial in.
    \item Let $\gamma_2(x)$ if $\exists y,z ~ sec(y,z) \land [y,x)\cong\gamma$ for any $\gamma\in\omega^\lambda$. Note that this definition is of complexity $\Sinf{\lambda+1}$ as $\lambda$ is a limit ordinal. It describes the points that are at the $\gamma$ position within a second copy of $\omega^{\lambda}$.
\end{enumerate}
Using these relations and those defined above, one can define the automorphism
orbits of tuples in $\B$.

In order to finish the proof of the claim, we appeal to Corollary \ref{lower} to
show that there is no $\Pinf{\lambda+3}$ Scott sentence.
This follows immediately as $\A\leq_{\lambda+1} \B$ is a direct consequence of $\L_1\leq_{\lambda+1}\L_2$.
\end{proof}

\begin{thebibliography}{9}
\bibitem{ash2000} C. Ash and J. Knight. \emph{Computable Structures and the Hyperarithmetical Hierarchy}. Studies in Logic and the Foundations of Mathematics144,2000.
\bibitem{AR20} R. Alvir and D. Rossegger. The complexity of Scott sentences of scattered linear orders. \emph{Journal of Symbolic Logic}85(3)(2020),1079--1101.
\bibitem{ash1986} C.J. Ash. Recursive labelling systems and stability of recursive structures in hyperarithmetical degrees. \emph{Transactions of the American Mathematical Society}298(2)(1986),497--514.
\bibitem{montalban2021} A. Montalban. \emph{Computable Structure Theory: Beyond the Arithmetic}. Draft,2021. \texttt{https://math.berkeley.edu/\string~antonio/CSTpart2\_DRAFT.pdf}.
\end{thebibliography}
\end{document}
